\section{Background}
\label{sec:background}

Research on political evasion and reply behaviour spans experimental and
observational traditions. Situational accounts of equivocation emphasise
avoid--avoid conflicts in which any direct answer is costly
\citep{bavelas_etal_1988}. Coding schemes for interviews and parliamentary
questions distinguish replies, intermediate forms, and non-replies, and have
been applied extensively to UK Prime Minister's Questions
\citep{bull_mayer_1993,bull_1994_identifying,bull_strawson_2020,bull_waddle_2019_jlac,bull_waddle_2023_jspp}.
Related work examines how issue pressure shapes equivocation in other
interview systems \citep{feldman_kinoshita_bull_2016}.

Spanish scholarship has analysed oral questions as a parliamentary genre and
as a site of politeness, image work, and institutional constraint
\citep{fuentes_2011_clr,fuentes_2012_ph,santos_lopez_2010_iberica}. Comparative
pragmatic work on evasive answering in French and Spanish political discourse
likewise treats responsiveness as a communicative achievement rather than a
mere topic match \citep{sanchez_villanueva_2022_thesis}. Computational shared
tasks on political question evasion are emerging
\citep{thomas_etal_2026_semeval}; the present study remains a human annotation
design and does not use model labels.

What remains underspecified for Spanish plenary control to the Prime Minister
is a census-style descriptive measure of first-response explicitness under a
fixed formation-vote baseline, with linkage between registered initiatives and
transcript turns made inspectable, and with question-side composition reported
before outcome interpretation. That measurement gap motivates the design below.
